\begin{figure}[htbp]
\centering
\begin{tikzpicture}[x=5.7mm,y=7mm,alt={An RGB565 halfword has bits fifteen through eleven for red, ten through five for green, and four through zero for blue. The example fields are red fifteen, green eleven and blue eight. Together they form hexadecimal 7968, stored as low byte 68 followed by high byte 79.}]
\fill[black!12] (0,0) rectangle (5,1);
\fill[black!22] (5,0) rectangle (11,1);
\fill[black!6] (11,0) rectangle (16,1);
\draw (0,0) rectangle (16,1) (5,0)--(5,1) (11,0)--(11,1);
\node at (2.5,.5) {Red: 15};\node at (8,.5) {Green: 11};\node at (13.5,.5) {Blue: 8};
\node[above] at (0,1) {15};\node[above] at (5,1) {10};\node[above] at (11,1) {4};\node[above] at (16,1) {0};
\node[below,align=center] at (8,-.25) {Numerical halfword: \texttt{7968}\\Stored bytes, increasing address: \texttt{68 79}};
\end{tikzpicture}
\caption{Field positions describe a numerical value. The byte order determines which part of that value is stored first. Shading distinguishes the fields without assigning display colors to their bits.}
\label{fig:revision-msa-fields}
\end{figure}
